\chapter{PHƯƠNG PHÁP}
\label{chap:methods}

\section{Tổng quan pipeline ba giai đoạn}
\label{sec:met-overview}

Hình~\ref{fig:pipeline} tóm tắt toàn bộ hệ thống. Giai đoạn A tiền huấn luyện
mô hình ngôn ngữ \textbf{ViGPT-nano} (101M tham số) trên 3,02 tỷ token tiếng
Việt. Giai đoạn B ghép bộ mã hóa thị giác SigLIP đóng băng qua projector
pixel-shuffle và tinh chỉnh có giám sát hai pha trên dữ liệu mô tả ảnh --
hỏi--đáp tiếng Việt. Giai đoạn C tinh chỉnh học tăng cường SCST với phần
thưởng CIDEr. Mỗi giai đoạn tạo một checkpoint độc lập
(\texttt{vigpt-base} $\to$ \texttt{vivlm-sft} $\to$ \texttt{vivlm-scst}),
cho phép đánh giá tách bạch đóng góp của từng bước.

\begin{figure}[H]
	\centering
	\begin{tikzpicture}[
		node distance=4mm and 6mm, font=\small,
		box/.style={draw, rounded corners=2pt, align=center, minimum height=9mm, inner sep=5pt},
		data/.style={box, fill=gray!10},
		model/.style={box, fill=blue!8},
		stage/.style={font=\bfseries\small}
	]
		\node[data] (fw) {FineWeb2-HQ \texttt{vie\_Latn}\\ 3{,}02 tỷ token};
		\node[model, right=of fw] (gpt) {ViGPT-nano 101M\\ (12 tầng Llama-style)};
		\node[data, right=of gpt] (sft) {UIT-ViIC + KTVIC\\ + OpenViVQA ($\approx$77k mẫu)};
		\node[model, right=of sft] (vlm) {ViVLM-nano\\ (+SigLIP, +projector)};
		\draw[-{Stealth}] (fw) -- node[above, stage] {A} (gpt);
		\draw[-{Stealth}] (gpt) -- node[above, stage] {B} (sft);
		\draw[-{Stealth}] (sft) -- node[above, stage] {C: SCST} (vlm);
	\end{tikzpicture}
	\caption{Pipeline ba giai đoạn: (A) tiền huấn luyện ngôn ngữ, (B) hợp nhất
	thị giác + SFT, (C) học tăng cường SCST.}
	\label{fig:pipeline}
\end{figure}

\section{Dữ liệu tiền huấn luyện}
\label{sec:met-pretrain-data}

Corpus lấy từ \textbf{FineWeb2-HQ} \autocite{messmer2025fineweb2hq} --- phân
tập tiếng Việt (\texttt{vie\_Latn}) đã lọc chất lượng bằng mô hình từ FineWeb2
\autocite{penedo2024fineweb}, tổng 63\,GB dạng parquet. Pipeline chuẩn bị dữ
liệu tải \emph{tuần tự từng file đến khi đủ chỉ tiêu token} (thay vì cố định
số file, vì mỗi file parquet chứa kèm cột embedding rất nặng khiến dung lượng
không tỷ lệ với lượng văn bản): 54 file đầu cho ra tập huấn luyện, file kế
tiếp --- không giao với tập huấn luyện --- cho tập kiểm định
(bảng~\ref{tab:pretrain-data}). Mỗi văn bản được mã hóa BPE rồi kết thúc bằng
token \texttt{<|endoftext|>}, ghi nối tiếp vào một file nhị phân
\texttt{uint16} (từ vựng 20.480 $<$ 65.536 nên 2 byte/token là đủ). Khi huấn
luyện, dữ liệu được đọc bằng \texttt{numpy.memmap} và lấy mẫu cửa sổ ngẫu
nhiên 1.024 token --- không nạp 6\,GB vào RAM, không nghẽn CPU.

\begin{table}[H]
	\centering
	\caption{Thống kê dữ liệu tiền huấn luyện.}
	\label{tab:pretrain-data}
	\begin{tabular}{lr}
		\hline
		Nguồn & FineWeb2-HQ \texttt{vie\_Latn} (54 + 1 file parquet) \\
		Token huấn luyện & 3.019.000.000 (6,0\,GB uint16) \\
		Token kiểm định & 30.600.000 (59\,MB, file riêng --- không rò rỉ) \\
		Cửa sổ ngữ cảnh & 1.024 token \\
		Phân cách văn bản & \texttt{<|endoftext|>} sau mỗi tài liệu \\
		\hline
	\end{tabular}
\end{table}

\section{Bộ tách từ}
\label{sec:met-tokenizer}

Bộ tách từ là BPE byte-level 20.480 mục, tự huấn luyện bằng thư viện
\texttt{tokenizers} trên $\approx$0,35\,GB văn bản trích từ file parquet đầu
tiên (đủ để ước lượng thống kê cặp byte của tiếngng Việt web). Bốn token đặc
biệt chiếm cố định các id 0--3: \texttt{<|endoftext|>} (kết thúc/phân cách),
\texttt{<|user|>}, \texttt{<|assistant|>} (định dạng hội thoại cho SFT) và
\texttt{<|image|>} (dự phòng). Kích thước 20.480 trùng với PhoGPT
\autocite{nguyen2023phogpt} để so sánh fertility trực tiếp
(mục~\ref{sec:exp-tokenizer}).

\section{Kiến trúc ViGPT-nano}
\label{sec:met-arch}

ViGPT-nano là transformer decoder-only Llama-style cài đặt từ đầu bằng
PyTorch (không dùng \texttt{nn.TransformerDecoder} hay thư viện
\texttt{transformers}), cấu hình ở bảng~\ref{tab:arch}. Mỗi khối:
pre-RMSNorm $\to$ self-attention đa đầu (RoPE trên $Q,K$, mặt nạ nhân quả,
FlashAttention) $\to$ cộng phần dư $\to$ pre-RMSNorm $\to$ SwiGLU $\to$ cộng
phần dư. Trọng số khởi tạo $\mathcal{N}(0, 0{,}02^2)$; riêng các lớp chiếu
cuối nhánh phần dư nhân thêm $1/\sqrt{2L}$ theo GPT-2
\autocite{radford2019gpt2} để phương sai tín hiệu không phình theo độ sâu.
Embedding trói trọng số với lớp đầu ra.

\begin{table}[H]
	\centering
	\caption{Cấu hình ViGPT-nano và phân bổ tham số.}
	\label{tab:arch}
	\begin{tabular}{lr|lr}
		\hline
		Số tầng $L$ & 12 & Embedding (trói với đầu ra) & 15,73M \\
		Chiều ẩn $d$ & 768 & Attention ($12 \times 2{,}36$M) & 28,3M \\
		Số đầu attention & 12 (mỗi đầu 64 chiều) & SwiGLU ($12 \times 4{,}72$M) & 56,6M \\
		Chiều ẩn SwiGLU & 2.048 & RMSNorm (25 lớp) & 0,02M \\
		Ngữ cảnh & 1.024 & & \\
		Từ vựng & 20.480 & \textbf{Tổng} & \textbf{100.682.496} \\
		\hline
	\end{tabular}
\end{table}

\section{Cấu hình tiền huấn luyện (giai đoạn A)}
\label{sec:met-pretrain}

Bảng~\ref{tab:pretrain-hp} liệt kê siêu tham số. Batch hiệu dụng 491.520
token/bước đạt bằng tích lũy gradient: 40 chuỗi $\times$ 1.024 token
$\times$ 12 lần tích lũy. 6.000 bước tương ứng $\approx$2,95 tỷ token
($\approx$0,98 epoch --- mỗi mẫu chỉ được thấy xấp xỉ một lần). Toàn bộ chạy
\texttt{bfloat16} (autocast) với \texttt{torch.compile}
\autocite{ansel2024pytorch2}.

\begin{table}[H]
	\centering
	\caption{Siêu tham số tiền huấn luyện.}
	\label{tab:pretrain-hp}
	\begin{tabular}{ll}
		\hline
		Optimizer & AdamW \autocite{loshchilov2019adamw}, $\beta = (0{,}9;\ 0{,}95)$, weight decay 0,1 \\
		 & (decay chỉ áp lên tham số $\ge$2 chiều; norm/bias miễn) \\
		Lịch học & warmup tuyến tính 500 bước $\to$ đỉnh $6 \cdot 10^{-4}$ $\to$ cosine về $6 \cdot 10^{-5}$ \\
		Batch & micro 40 $\times$ ctx 1.024 $\times$ accum 12 = 491.520 token/bước \\
		Số bước & 6.000 ($\approx$2,95 tỷ token $\approx$ 0,98 epoch) \\
		Gradient clipping & chuẩn toàn cục 1,0 \\
		Độ chính xác & bfloat16 autocast + FlashAttention + \texttt{torch.compile} \\
		Đánh giá / checkpoint & val mỗi 250 bước (50 batch cố định); checkpoint mỗi 500 bước \\
		\hline
	\end{tabular}
\end{table}

Vì phiên huấn luyện kéo dài nhiều giờ, tính \textbf{khôi phục chính xác} được
thiết kế và kiểm thử tường minh: checkpoint lưu đủ trọng số, trạng thái
optimizer, bước hiện tại và \emph{mọi trạng thái sinh số ngẫu nhiên} (bộ sinh
lấy mẫu dữ liệu và RNG của PyTorch trên mọi thiết bị). Unit test khóa bất biến
``huấn luyện 6 bước liền mạch $=$ huấn luyện 3 bước + khôi phục + 3 bước''
chính xác đến từng bit trọng số. Một lỗi thực tế đáng ghi nhận: tensor trạng
thái RNG bị \texttt{map\_location} chuyển sang thiết bị tăng tốc khi nạp
checkpoint, trong khi API khôi phục đòi hỏi tensor trên CPU --- được vá và
khóa bằng test riêng.

\section{Hợp nhất thị giác và SFT (giai đoạn B)}
\label{sec:met-fusion}

\subsection{Kiến trúc hợp nhất}

Ảnh $224 \times 224$ qua SigLIP-B/16 \emph{đóng băng} cho 196 patch 768
chiều. Pixel-shuffle gộp $2 \times 2$ patch thành một token 3.072 chiều
($196 \to 49$ token), rồi projector hai lớp
($3.072 \to 768$, SiLU, $768 \to 768$) chiếu vào không gian embedding của
ViGPT-nano. Chuỗi đầu vào cuối cùng là
$[\,49 \text{ token ảnh};\ \text{token văn bản}\,]$ --- các tầng transformer
sẵn có đảm nhiệm hợp nhất, đúng tinh thần LLaVA (mục~\ref{sec:bg-fusion}).

\begin{figure}[H]
	\centering
	\begin{tikzpicture}[
		node distance=3mm and 5mm, font=\small,
		box/.style={draw, rounded corners=2pt, align=center, minimum height=8mm, inner sep=4pt},
		frozen/.style={box, fill=cyan!8},
		train/.style={box, fill=orange!12}
	]
		\node[frozen] (sig) {SigLIP-B/16 (đóng băng)\\ ảnh $\to 196 \times 768$};
		\node[train, right=of sig] (proj) {pixel-shuffle $196{\to}49$\\ + MLP $3072{\to}768$};
		\node[train, right=of proj] (gpt) {ViGPT-nano\\ 12 tầng};
		\node[below=6mm of gpt, box, fill=gray!10] (txt) {\texttt{<|user|>} Ảnh có gì? \texttt{<|assistant|>} \ldots};
		\draw[-{Stealth}] (sig) -- (proj);
		\draw[-{Stealth}] (proj) -- node[above]{ghép trước} (gpt);
		\draw[-{Stealth}] (txt) -- (gpt);
	\end{tikzpicture}
	\caption{Hợp nhất kiểu LLaVA: token ảnh (nén 4 lần bằng pixel-shuffle)
	ghép trước chuỗi token văn bản; chỉ projector và ViGPT-nano được huấn luyện.}
	\label{fig:fusion}
\end{figure}

\subsection{Dữ liệu SFT}

Ba nguồn tiếng Việt được chuẩn hóa về một định dạng thống nhất
(ảnh, câu lệnh, câu trả lời), tổng $\approx$77 nghìn mẫu
(bảng~\ref{tab:sft-data}); 2\% tách làm tập kiểm định theo băm định danh
\emph{ảnh} (mọi mẫu của cùng một ảnh nằm cùng phía, tránh rò rỉ). Với dữ liệu
mô tả ảnh, câu lệnh được đa dạng hóa từ 5 mẫu câu (``Mô tả bức ảnh.'', ``Ảnh
này có gì?''\ldots) chọn tất định theo băm, tránh mô hình học vẹt một câu lệnh
duy nhất. Chuỗi huấn luyện theo định dạng
\texttt{<|user|>}~$q$~\texttt{<|assistant|>}~$a$~\texttt{<|endoftext|>}; nhãn
bị che ($-100$) trên toàn bộ 49 token ảnh và phần câu lệnh --- gradient chỉ
đến từ câu trả lời (mục~\ref{sec:bg-sft}).

\begin{table}[H]
	\centering
	\caption{Ba nguồn dữ liệu SFT tiếng Việt.}
	\label{tab:sft-data}
	\begin{tabularx}{\textwidth}{lXrr}
		\hline
		\textbf{Bộ} & \textbf{Miền} & \textbf{Ảnh} & \textbf{Mẫu} \\
		\hline
		UIT-ViIC \autocite{lam2020viic} & thể thao (ảnh COCO), mô tả & $\approx$3,8k & $\approx$19k \\
		KTVIC \autocite{pham2024ktvic} & đời sống hằng ngày, mô tả & $\approx$4,3k & $\approx$21k \\
		OpenViVQA \autocite{nguyen2023openvivqa} & hỏi--đáp mở (VLSP 2023) & $\approx$11k & $\approx$37k \\
		\hline
	\end{tabularx}
\end{table}

\subsection{Huấn luyện hai pha}

Pha 1 (\emph{làm nóng projector}): đóng băng toàn bộ ViGPT-nano, chỉ huấn
luyện projector 500 bước với learning rate $10^{-3}$ --- cho projector ``bắt
sóng'' không gian embedding trước, tránh gradient hỗn loạn từ một projector
khởi tạo ngẫu nhiên phá hỏng phần ngôn ngữ đã pretrain. Pha 2 (\emph{toàn
phần}): mở khóa ViGPT-nano + projector (SigLIP vẫn đóng băng), 2 epoch với
learning rate $10^{-4}$ cosine về $10^{-5}$, micro-batch 24, tích lũy 4,
độ dài văn bản tối đa 256 token.

\section{Học tăng cường SCST (giai đoạn C)}
\label{sec:met-scst}

Từ checkpoint SFT, mỗi bước SCST lấy 32 ảnh từ nguồn dữ liệu mô tả (kèm các
câu tham chiếu), giải mã hai lần dưới \texttt{no\_grad}: tham lam
($\hat{w}$, baseline) và lấy mẫu ($w^s$, nhiệt độ 1,0); phần thưởng là CIDEr
từng ảnh. Log-xác suất của chuỗi lấy mẫu sau đó được tính lại bằng \emph{một}
lượt truyền teacher-forcing có gradient (thay vì giữ đồ thị tính toán qua 40
bước giải mã --- tiết kiệm bộ nhớ tương đương một bước huấn luyện thường), và
loss là
\begin{equation}
	\mathcal{L}_{\mathrm{SCST}} = -\bigl(r(w^s) - r(\hat{w})\bigr) \sum_{t} \log p_\theta(w^s_t \mid w^s_{<t}, I),
	\label{eq:scst-loss}
\end{equation}
với mặt nạ cắt tại \texttt{<|endoftext|>} đầu tiên của từng chuỗi trong
batch. Chạy 400 bước, learning rate $10^{-5}$, tối đa 40 token sinh mới.

\section{Cài đặt và kiểm thử}
\label{sec:met-infra}

Toàn bộ phần học được (khối transformer, projector, ba vòng lặp huấn luyện,
bộ tách từ) được cài đặt từ đầu bằng PyTorch, không dùng \texttt{transformers}
của Hugging Face cho phần ngôn ngữ. Pipeline được phát triển kiểu test-driven:
\textbf{52 unit test} chạy ngoại tuyến trên máy cá nhân chỉ dùng CPU (mô hình
thu nhỏ, bộ mã hóa thị giác giả lập, dữ liệu tự sinh), khóa các bất biến quan
trọng --- RoPE bảo toàn chuẩn, mặt nạ nhân quả, đếm đúng 100.682.496 tham số,
trói trọng số, che nhãn SFT đúng vị trí, khôi phục checkpoint chính xác từng
bit, căn chỉnh log-xác suất SCST qua 49 token ảnh, và một test hồi quy cho một
lỗi đường dẫn ảnh phát hiện khi rà soát chéo. Khi đánh giá sinh chuỗi bằng
beam search, điểm mỗi beam được chuẩn hóa theo độ dài (mũ 0,7) để tránh thiên
vị câu ngắn.
